% THOTH CLOUD - Implementation Report
% INF4107 - TP4: Mini Cloud PaaS
% 2025-2026

\documentclass[12pt,a4paper]{article}

% Packages
\usepackage[utf8]{inputenc}
\usepackage[french]{babel}
\usepackage[T1]{fontenc}
\usepackage{geometry}
\usepackage{graphicx}
\usepackage{xcolor}
\usepackage{listings}
\usepackage{tcolorbox}
\usepackage{fancyhdr}
\usepackage{hyperref}
\usepackage{titlesec}
\usepackage{enumitem}
\usepackage{tikz}
\usepackage{pgfplots}
\usepackage{tabularx}
\usepackage{booktabs}
\usepackage{float}

% Geometry
\geometry{margin=2.5cm}

% Color definitions - Green-Blue theme
\definecolor{primaryBlue}{RGB}{41, 128, 185}
\definecolor{secondaryGreen}{RGB}{39, 174, 96}
\definecolor{accentTeal}{RGB}{26, 188, 156}
\definecolor{darkBlue}{RGB}{44, 62, 80}
\definecolor{lightBlue}{RGB}{174, 214, 241}
\definecolor{lightGreen}{RGB}{163, 228, 215}
\definecolor{codeBackground}{RGB}{236, 240, 241}

% Hyperref setup
\hypersetup{
    colorlinks=true,
    linkcolor=primaryBlue,
    filecolor=secondaryGreen,
    urlcolor=accentTeal,
    citecolor=primaryBlue,
    pdftitle={THOTH CLOUD - Rapport d'Implémentation},
    pdfauthor={Équipe THOTH CLOUD}
}

% Title formatting
\titleformat{\section}
{\color{primaryBlue}\normalfont\Large\bfseries}
{\color{primaryBlue}\thesection}{1em}{}

\titleformat{\subsection}
{\color{secondaryGreen}\normalfont\large\bfseries}
{\color{secondaryGreen}\thesubsection}{1em}{}

\titleformat{\subsubsection}
{\color{accentTeal}\normalfont\normalsize\bfseries}
{\color{accentTeal}\thesubsubsection}{1em}{}

% Custom tcolorbox styles
\tcbuselibrary{skins,breakable}

\newtcolorbox{infobox}[1]{
    colback=lightBlue,
    colframe=primaryBlue,
    fonttitle=\bfseries,
    title=#1,
    breakable,
    enhanced
}

\newtcolorbox{successbox}[1]{
    colback=lightGreen,
    colframe=secondaryGreen,
    fonttitle=\bfseries,
    title=#1,
    breakable,
    enhanced
}

\newtcolorbox{codebox}[1]{
    colback=codeBackground,
    colframe=darkBlue,
    fonttitle=\bfseries\ttfamily,
    title=#1,
    breakable,
    enhanced,
    listing only
}

% Listings configuration
\lstset{
    basicstyle=\small\ttfamily,
    backgroundcolor=\color{codeBackground},
    keywordstyle=\color{primaryBlue}\bfseries,
    commentstyle=\color{secondaryGreen}\itshape,
    stringstyle=\color{accentTeal},
    numbers=left,
    numberstyle=\tiny\color{darkBlue},
    stepnumber=1,
    numbersep=8pt,
    showstringspaces=false,
    breaklines=true,
    frame=single,
    rulecolor=\color{darkBlue},
    tabsize=4,
    captionpos=b
}

% Header and footer
\pagestyle{fancy}
\fancyhf{}
\fancyhead[L]{\color{primaryBlue}\textbf{THOTH CLOUD}}
\fancyhead[R]{\color{secondaryGreen}Rapport d'Implémentation}
\fancyfoot[C]{\color{darkBlue}\thepage}
\renewcommand{\headrulewidth}{2pt}
\renewcommand{\headrule}{\hbox to\headwidth{\color{accentTeal}\leaders\hrule height \headrulewidth\hfill}}

% Document
\begin{document}

% Title page
\begin{titlepage}
    \centering
    \vspace*{2cm}
    
    {\Huge\bfseries\color{primaryBlue} THOTH CLOUD\par}
    \vspace{0.5cm}
    {\Large\color{secondaryGreen} Plateforme Cloud Multi-Services\par}
    \vspace{2cm}
    
    {\LARGE\bfseries\color{darkBlue} Rapport d'Implémentation Backend\par}
    \vspace{0.5cm}
    {\large\color{accentTeal} TP4 : Évolution vers un Mini Cloud PaaS\par}
    \vspace{3cm}
    
    \begin{tikzpicture}
        \draw[fill=lightBlue, draw=primaryBlue, line width=2pt] (0,0) rectangle (10,3);
        \node[align=center] at (5,2.2) {\Large\bfseries\color{darkBlue} INF4107};
        \node[align=center] at (5,1.5) {\large\color{primaryBlue} Virtualisation et Cloud Computing};
        \node[align=center] at (5,0.8) {\color{secondaryGreen} Master 1 - Systèmes \& Réseaux};
    \end{tikzpicture}
    
    \vfill
    
    {\large
    \begin{tabular}{ll}
        \textbf{\color{primaryBlue}Enseignants:} & \color{darkBlue}Valery MONTHE \& Gildas NGOUANFO \\
        \textbf{\color{primaryBlue}Année:} & \color{darkBlue}2025-2026 \\
    \end{tabular}
    }
    
    \vspace{1cm}
    {\color{accentTeal}\rule{0.8\textwidth}{2pt}}
\end{titlepage}

% Table of contents
\tableofcontents
\newpage

% Introduction
\section{Introduction}

\begin{infobox}{Contexte du Projet}
THOTH CLOUD est une plateforme cloud complète implémentée en C++ qui fournit des services IaaS, PaaS et SaaS. Ce rapport documente l'architecture backend et les implémentations techniques réalisées conformément aux exigences du TP4.
\end{infobox}

\subsection{Objectifs du Projet}

Le projet vise à créer un cloud multi-tenant avec:
\begin{itemize}[leftmargin=*, itemsep=0.3em]
    \item \textbf{IaaS amélioré}: Multi-utilisateurs, isolation réseau, gestion multi-hôtes
    \item \textbf{PaaS}: Déploiement automatique de clusters Docker Swarm et applications
    \item \textbf{SaaS}: Plateforme collaborative (OwnCloud + OnlyOffice)
    \item \textbf{Facturation}: Système de billing basé sur la consommation
\end{itemize}

\subsection{Lien GitHub}

\begin{successbox}{Dépôt du Projet}
\centering
\url{https://github.com/votre-organisation/thoth-cloud}
\end{successbox}

% Architecture
\section{Architecture Globale}

\subsection{Vue d'Ensemble}

\begin{figure}[H]
\centering
\begin{tikzpicture}[scale=0.9, every node/.style={scale=0.9}]
    % Frontend Layer
    \draw[fill=lightBlue, draw=primaryBlue, line width=1.5pt] (0,8) rectangle (12,10);
    \node[align=center] at (6,9) {\Large\bfseries\color{darkBlue} Frontend Web Interface};
    \node[align=center, font=\small] at (6,8.4) {\color{primaryBlue} React/Vue.js - Interface Utilisateur};
    
    % Backend Layer
    \draw[fill=lightGreen, draw=secondaryGreen, line width=1.5pt] (0,5.5) rectangle (12,7.5);
    \node[align=center] at (6,6.8) {\Large\bfseries\color{darkBlue} Backend API (C++)};
    \node[align=center, font=\small] at (6,6.2) {\color{secondaryGreen} REST API - Gestion des Ressources};
    
    % Core Services
    \draw[fill=codeBackground, draw=accentTeal, line width=1pt] (0.5,3.5) rectangle (3.5,5);
    \node[align=center, font=\small\bfseries] at (2,4.5) {\color{primaryBlue} VM\\Operations};
    
    \draw[fill=codeBackground, draw=accentTeal, line width=1pt] (4,3.5) rectangle (7,5);
    \node[align=center, font=\small\bfseries] at (5.5,4.5) {\color{primaryBlue} Host\\Manager};
    
    \draw[fill=codeBackground, draw=accentTeal, line width=1pt] (7.5,3.5) rectangle (10.5,5);
    \node[align=center, font=\small\bfseries] at (9,4.5) {\color{primaryBlue} Network\\Manager};
    
    % Infrastructure Layer
    \draw[fill=lightBlue, draw=primaryBlue, line width=1.5pt] (0,0.5) rectangle (5.5,3);
    \node[align=center] at (2.75,2.2) {\large\bfseries\color{darkBlue} KVM Hosts};
    \node[align=center, font=\small] at (2.75,1.5) {\color{primaryBlue} Libvirt + QEMU};
    
    \draw[fill=lightGreen, draw=secondaryGreen, line width=1.5pt] (6.5,0.5) rectangle (12,3);
    \node[align=center] at (9.25,2.2) {\large\bfseries\color{darkBlue} Docker Swarm};
    \node[align=center, font=\small] at (9.25,1.5) {\color{secondaryGreen} Container Orchestration};
    
    % Arrows
    \draw[->, line width=1.5pt, color=accentTeal] (6,8) -- (6,7.5);
    \draw[->, line width=1pt, color=primaryBlue] (2,5.5) -- (2,3);
    \draw[->, line width=1pt, color=primaryBlue] (5.5,5.5) -- (5.5,3);
    \draw[->, line width=1pt, color=primaryBlue] (9,5.5) -- (9,3);
\end{tikzpicture}
\caption{Architecture en Couches de THOTH CLOUD}
\end{figure}

\subsection{Stack Technologique}

\begin{table}[H]
\centering
\begin{tabular}{@{}ll@{}}
\toprule
\textbf{\color{primaryBlue}Composant} & \textbf{\color{secondaryGreen}Technologie} \\
\midrule
Langage Backend & C++17 \\
Framework HTTP & cpp-httplib \\
Sérialisation JSON & nlohmann/json \\
Virtualisation & Libvirt + QEMU/KVM \\
Conteneurisation & Docker + Docker Swarm \\
Orchestration & Cloud-init, Ansible \\
Base de données & PostgreSQL, MariaDB \\
Réseau & Libvirt networks, iptables \\
\bottomrule
\end{tabular}
\caption{Technologies Utilisées}
\end{table}

% Partie 1
\section{Partie 1: Amélioration du Service IaaS}

\subsection{Multi-Tenant avec Isolation Utilisateur}

\begin{successbox}{Isolation des Ressources}
Chaque utilisateur dispose de ressources isolées et ne peut accéder qu'à ses propres VMs, réseaux et clusters.
\end{successbox}

\subsubsection{Système de Nommage des VMs}

Implémentation dans \texttt{utils.hpp} - classe \texttt{VMNameManager}:

\begin{lstlisting}[language=C++, caption=Génération de noms de VMs uniques]
std::string createVMName(const std::string& username, 
                        const std::string& userVMName) {
    std::string cleanUser = sanitize(username);
    std::string cleanVMName = sanitize(userVMName);
    std::string timestamp = getTimestamp();
    
    // Format: username__vmname__timestamp
    return cleanUser + SEPARATOR + cleanVMName + SEPARATOR + timestamp;
}
\end{lstlisting}

\textbf{Avantages:}
\begin{itemize}
    \item Noms globalement uniques
    \item Identification rapide du propriétaire
    \item Impossible pour un utilisateur de deviner les noms des VMs d'autrui
\end{itemize}

\subsubsection{Filtrage par Utilisateur}

\begin{lstlisting}[language=C++, caption=Filtrage des VMs par utilisateur]
int UserOperations::listUserDomains(virDomainPtr **domains, 
                                    std::string username, 
                                    int flags) {
    virDomainPtr* allDomains = nullptr;
    int numDomains = virConnectListAllDomains(conn, &allDomains, flags);
    
    // Filter by username prefix
    VMNameManager::filterUserVMs(allDomains, username, 
                                domains, numDomains, 
                                &numUserDomains);
    return numUserDomains;
}
\end{lstlisting}

\subsection{Découplage de l'Application}

\begin{infobox}{Architecture Découplée}
L'application backend est séparée des hôtes KVM et communique via SSH/libvirt remote.
\end{infobox}

\subsubsection{Remote Executor}

Classe \texttt{RemoteExecutor} (\texttt{remote\_executor.hpp/cpp}):

\begin{lstlisting}[language=C++, caption=Exécution de commandes à distance]
class RemoteExecutor {
private:
    virConnectPtr conn;
    bool isRemote;
    std::string remoteUser;
    std::string remoteHost;
    std::string sshKeyFile;
    
    std::string buildSSHCommand(const std::string& command) const {
        if (!isRemote) return command;
        
        std::stringstream ssh;
        ssh << "ssh ";
        if (!sshKeyFile.empty()) {
            ssh << "-i " << sshKeyFile << " ";
        }
        ssh << "-o StrictHostKeyChecking=no ";
        ssh << remoteUser << "@" << remoteHost << " ";
        ssh << "'" << command << "'";
        return ssh.str();
    }
    
public:
    ExecResult execute(const std::string& command) const;
    bool fileExists(const std::string& path) const;
    long long getAvailableDiskSpace(const std::string& path) const;
};
\end{lstlisting}

\textbf{Fonctionnalités:}
\begin{itemize}
    \item Détection automatique du mode (local/remote)
    \item Authentification par clé SSH
    \item Vérification de fichiers distants
    \item Gestion de l'espace disque
\end{itemize}

\subsection{Isolation Réseau}

\begin{successbox}{Stratégie d'Isolation}
\begin{itemize}
    \item Un réseau libvirt dédié par utilisateur
    \item Un réseau IP séparé par cluster Docker Swarm
    \item Attribution automatique de sous-réseaux
\end{itemize}
\end{successbox}

\subsubsection{Network Manager}

Implémentation dans \texttt{network\_manager.hpp/cpp}:

\begin{lstlisting}[language=C++, caption=Création de réseau utilisateur]
json NetworkManager::createUserNetwork(const std::string& username, 
                                      const std::string& networkName) {
    json result;
    
    // Check user network limit
    if (getUserNetworkCount(username) >= MAX_USER_NETWORKS) {
        result["error"] = "User network limit reached";
        return result;
    }
    
    // Generate unique subnet
    std::string subnet = generateSubnet();
    std::string networkId = generateNetworkId();
    
    // Create libvirt network XML
    std::string xml = generateUserNetworkXML(networkId, 
                                            displayName, 
                                            subnet);
    
    // Define and start network
    virNetworkPtr network = virNetworkDefineXML(conn, xml.c_str());
    virNetworkCreate(network);
    virNetworkSetAutostart(network, 1);
    
    // Save configuration
    saveNetworksConfig();
    
    result["success"] = true;
    result["networkId"] = networkId;
    return result;
}
\end{lstlisting}

\textbf{Génération XML:}
\begin{lstlisting}[language=C++, caption=XML pour réseau utilisateur]
std::string NetworkManager::generateUserNetworkXML(
    const std::string& networkId, 
    const std::string& displayName,
    const std::string& subnet) {
    
    std::stringstream xml;
    xml << "<network>"
        << "<name>" << networkId << "</name>"
        << "<bridge name='virbr-" << networkId.substr(0,8) << "'/>"
        << "<forward mode='nat'/>"
        << "<ip address='" << subnet << ".1' netmask='255.255.255.0'>"
        << "<dhcp>"
        << "<range start='" << subnet << ".100' "
        << "end='" << subnet << ".200'/>"
        << "</dhcp>"
        << "</ip>"
        << "</network>";
    return xml.str();
}
\end{lstlisting}

\subsection{Flavors et Gestion des Ressources}

\begin{infobox}{Système de Flavors}
Les flavors définissent des configurations prédéfinies de ressources avec tarification associée.
\end{infobox}

\subsubsection{Définition des Flavors}

Structure dans \texttt{flavor\_manager.hpp}:

\begin{lstlisting}[language=C++, caption=Structure Flavor]
struct Flavor {
    std::string id;
    std::string name;
    std::string description;
    int vcpus;      // Nombre de vCPUs
    int ram;        // RAM en MB
    int disk;       // Disque en GB
    int price;      // Prix en FCFA/mois
    bool active;
};
\end{lstlisting}

\subsubsection{Flavors par Défaut}

\begin{table}[H]
\centering
\begin{tabularx}{\textwidth}{@{}lXcccc@{}}
\toprule
\textbf{ID} & \textbf{Description} & \textbf{vCPU} & \textbf{RAM} & \textbf{Disk} & \textbf{Prix} \\
\midrule
nano & Entry level - Testing & 1 & 512 MB & 5 GB & 800 FCFA \\
micro & Small apps & 1 & 1 GB & 10 GB & 1500 FCFA \\
small & Light workloads & 1 & 2 GB & 15 GB & 3500 FCFA \\
medium & Standard workloads & 2 & 4 GB & 20 GB & 6500 FCFA \\
large & Heavy workloads & 2 & 8 GB & 40 GB & 12000 FCFA \\
xlarge & Database \& Analytics & 4 & 16 GB & 80 GB & 22000 FCFA \\
\bottomrule
\end{tabularx}
\caption{Catalogue de Flavors}
\end{table}

\subsubsection{Validation des Ressources}

\begin{lstlisting}[language=C++, caption=Vérification de disponibilité]
json HostManager::checkResourceAvailability(int memory, 
                                           int cpu, 
                                           long long disk) {
    json result;
    result["available"] = false;
    
    for (auto& [id, host] : hosts) {
        if (!host.active) continue;
        
        updateHostResources(host);
        
        bool canAccommodate = 
            host.availableMemory >= (memory * 1024) &&
            host.availableCPUs >= cpu &&
            host.availableDisk >= disk;
        
        if (canAccommodate) {
            result["available"] = true;
            result["selectedHost"] = id;
            break;
        }
    }
    
    return result;
}
\end{lstlisting}

\subsection{Gestion Multi-Hôtes}

\begin{successbox}{Pool d'Hôtes KVM}
Le système gère plusieurs hôtes KVM avec sélection automatique basée sur différentes stratégies.
\end{successbox}

\subsubsection{Stratégies de Sélection d'Hôte}

\textbf{1. Least Used (Moins Chargé)} - Stratégie par défaut

\begin{lstlisting}[language=C++, caption=Stratégie Least Used]
std::string HostManager::selectByLeastUsed(int mem, int cpu, 
                                          long long disk) {
    std::string bestHost;
    double bestScore = -1;
    
    for (auto& [id, host] : hosts) {
        if (!host.active) continue;
        
        updateHostResources(host);
        
        // Check sufficient resources
        if (host.availableMemory < mem * 1024 ||
            host.availableCPUs < cpu ||
            host.availableDisk < disk) {
            continue;
        }
        
        // Score based on free resources percentage
        double memScore = (host.availableMemory / 
                          (double)host.totalMemory) * 100;
        double cpuScore = (host.availableCPUs / 
                          (double)host.totalCPUs) * 100;
        double diskScore = (host.availableDisk / 
                           (double)host.totalDisk) * 100;
        
        // Weighted: 40% mem, 40% CPU, 20% disk
        double score = memScore * 0.4 + cpuScore * 0.4 + 
                      diskScore * 0.2;
        
        if (score > bestScore) {
            bestScore = score;
            bestHost = id;
        }
    }
    
    return bestHost;
}
\end{lstlisting}

\textbf{2. Round Robin}

\begin{lstlisting}[language=C++, caption=Stratégie Round Robin]
std::string HostManager::selectByRoundRobin(int mem, int cpu, 
                                           long long disk) {
    if (hosts.empty()) return "";
    
    std::vector<std::string> hostIds;
    for (const auto& [id, host] : hosts) {
        if (host.active) hostIds.push_back(id);
    }
    
    size_t attempts = 0;
    while (attempts < hostIds.size()) {
        size_t currentIndex = (roundRobinIndex + attempts) % 
                             hostIds.size();
        std::string hostId = hostIds[currentIndex];
        auto& host = hosts[hostId];
        
        updateHostResources(host);
        
        if (host.availableMemory >= mem * 1024 &&
            host.availableCPUs >= cpu &&
            host.availableDisk >= disk) {
            
            roundRobinIndex = (currentIndex + 1) % hostIds.size();
            return hostId;
        }
        attempts++;
    }
    
    return "";
}
\end{lstlisting}

\textbf{3. Best Fit}

\begin{lstlisting}[language=C++, caption=Stratégie Best Fit]
std::string HostManager::selectByBestFit(int mem, int cpu, 
                                        long long disk) {
    std::string bestHost;
    double bestFit = std::numeric_limits<double>::max();
    
    for (auto& [id, host] : hosts) {
        if (!host.active) continue;
        
        updateHostResources(host);
        
        if (host.availableMemory < mem * 1024 ||
            host.availableCPUs < cpu ||
            host.availableDisk < disk) {
            continue;
        }
        
        // Calculate "waste" - leftover resources
        double memWaste = (host.availableMemory / 1024.0 - mem) / 
                         (double)host.totalMemory;
        double cpuWaste = (host.availableCPUs - cpu) / 
                         (double)host.totalCPUs;
        double diskWaste = (host.availableDisk - disk) / 
                          (double)host.totalDisk;
        
        double wasteScore = memWaste * 0.4 + cpuWaste * 0.4 + 
                           diskWaste * 0.2;
        
        if (wasteScore < bestFit) {
            bestFit = wasteScore;
            bestHost = id;
        }
    }
    
    return bestHost;
}
\end{lstlisting}

\subsubsection{Autres Stratégies Possibles}

\begin{enumerate}
    \item \textbf{Priority-Based}: Hôtes avec priorité différente selon leur capacité
    \item \textbf{Geographic}: Sélection basée sur la proximité géographique
    \item \textbf{Cost-Optimized}: Minimisation des coûts opérationnels
    \item \textbf{Power-Aware}: Réduction de la consommation électrique
    \item \textbf{Load-Balanced}: Équilibrage de charge avec seuils
\end{enumerate}

% Partie 2
\section{Partie 2: Cluster Docker Swarm}

\subsection{Déploiement Automatique}

\begin{infobox}{Processus de Création}
Le système déploie automatiquement un cluster Swarm avec le nombre souhaité de managers et workers.
\end{infobox}

\subsubsection{Architecture du Cluster}

\begin{lstlisting}[language=C++, caption=Création de cluster]
json SwarmOperations::createSwarmCluster(
    const std::string& clusterName,
    const std::string& owner,
    int numManagers,
    int numWorkers) {
    
    json result;
    
    // 1. Create dedicated network for cluster
    json networkResult = networkMgr->createSwarmNetwork(
        clusterName, owner);
    
    if (!networkResult["success"].get<bool>()) {
        result["error"] = "Failed to create cluster network";
        return result;
    }
    
    std::string networkId = networkResult["network"]["networkId"];
    
    // 2. Deploy manager nodes
    for (int i = 0; i < numManagers; i++) {
        json vmConfig = {
            {"hostname", clusterName + "-manager-" + 
                        std::to_string(i+1)},
            {"owner", owner},
            {"role", "manager"},
            {"network", networkId},
            {"flavor", "medium"},
            {"baseImage", "ubuntu-22.04"}
        };
        
        json vmResult = deploySwarmNode(vmConfig);
        if (!vmResult["success"].get<bool>()) {
            // Rollback on failure
            deleteCluster(clusterId);
            return vmResult;
        }
    }
    
    // 3. Deploy worker nodes
    for (int i = 0; i < numWorkers; i++) {
        json vmConfig = {
            {"hostname", clusterName + "-worker-" + 
                        std::to_string(i+1)},
            {"owner", owner},
            {"role", "worker"},
            {"network", networkId},
            {"flavor", "small"},
            {"baseImage", "ubuntu-22.04"}
        };
        
        deploySwarmNode(vmConfig);
    }
    
    // 4. Initialize Swarm on first manager
    initializeSwarm(clusterName);
    
    result["success"] = true;
    result["clusterId"] = clusterId;
    return result;
}
\end{lstlisting}

\subsubsection{Installation Docker via Cloud-Init}

\begin{lstlisting}[language=bash, caption=Cloud-init pour Docker]
#cloud-config
package_update: true
package_upgrade: true

packages:
  - docker.io
  - docker-compose

runcmd:
  # Enable and start Docker
  - systemctl enable docker
  - systemctl start docker
  
  # Add user to docker group
  - usermod -aG docker ubuntu
  
  # Open Swarm ports
  - ufw allow 2377/tcp   # Cluster management
  - ufw allow 7946/tcp   # Node communication
  - ufw allow 7946/udp
  - ufw allow 4789/udp   # Overlay network
  
  # Install Docker Compose
  - curl -L "https://github.com/docker/compose/releases/latest/download/docker-compose-$(uname -s)-$(uname -m)" -o /usr/local/bin/docker-compose
  - chmod +x /usr/local/bin/docker-compose
\end{lstlisting}

\subsubsection{Initialisation du Swarm}

\begin{lstlisting}[language=C++, caption=Initialisation Swarm Manager]
json SwarmOperations::initializeSwarmOnManager(
    const std::string& vmName,
    const std::string& managerIP) {
    
    json result;
    
    // Initialize swarm
    std::string initCmd = "docker swarm init --advertise-addr " + 
                         managerIP;
    auto initResult = remoteExec->execute(initCmd);
    
    if (!initResult.success()) {
        result["error"] = "Failed to initialize swarm";
        return result;
    }
    
    // Get manager token
    std::string managerTokenCmd = 
        "docker swarm join-token manager -q";
    auto tokenResult = remoteExec->execute(managerTokenCmd);
    std::string managerToken = tokenResult.output;
    
    // Get worker token
    std::string workerTokenCmd = 
        "docker swarm join-token worker -q";
    auto workerTokenResult = remoteExec->execute(workerTokenCmd);
    std::string workerToken = workerTokenResult.output;
    
    result["success"] = true;
    result["managerToken"] = managerToken;
    result["workerToken"] = workerToken;
    
    return result;
}
\end{lstlisting}

\subsubsection{Jonction des Workers}

\begin{lstlisting}[language=C++, caption=Ajout de worker au cluster]
json SwarmOperations::joinWorkerToSwarm(
    const std::string& vmName,
    const std::string& workerIP,
    const std::string& managerIP,
    const std::string& token) {
    
    std::string joinCmd = 
        "docker swarm join --token " + token + 
        " " + managerIP + ":2377";
    
    auto result = remoteExec->execute(joinCmd);
    
    json response;
    response["success"] = result.success();
    
    return response;
}
\end{lstlisting}

\subsection{Exemple de Déploiement}

\begin{successbox}{Cluster 1 Manager + 2 Workers}
Configuration recommandée pour démonstration
\end{successbox}

\begin{table}[H]
\centering
\begin{tabular}{@{}llll@{}}
\toprule
\textbf{Nœud} & \textbf{Rôle} & \textbf{Flavor} & \textbf{IP} \\
\midrule
demo-manager-1 & Manager & medium (2 vCPU, 4GB) & 192.168.100.101 \\
demo-worker-1 & Worker & small (1 vCPU, 2GB) & 192.168.100.102 \\
demo-worker-2 & Worker & small (1 vCPU, 2GB) & 192.168.100.103 \\
\bottomrule
\end{tabular}
\caption{Configuration du Cluster de Démonstration}
\end{table}

% Partie 3
\section{Partie 3: Service PaaS}

\subsection{Catalogue d'Applications}

\begin{infobox}{Applications Disponibles}
Déploiement en un clic de 8 applications populaires avec bases de données mutualisées.
\end{infobox}

\subsubsection{Liste des Applications}

\begin{table}[H]
\centering
\begin{tabularx}{\textwidth}{@{}lXl@{}}
\toprule
\textbf{Application} & \textbf{Description} & \textbf{Base} \\
\midrule
WordPress & CMS et blogging & MariaDB \\
Odoo & ERP/CRM & PostgreSQL \\
Mattermost & Communication d'équipe & PostgreSQL \\
OwnCloud & Stockage cloud & MariaDB \\
OnlyOffice & Suite bureautique & PostgreSQL \\
OpenLDAP & Annuaire LDAP & - \\
Moodle & Plateforme e-learning & MariaDB \\
PrestaShop & E-commerce & MariaDB \\
\bottomrule
\end{tabularx}
\caption{Catalogue d'Applications PaaS}
\end{table}

\subsection{Mutualisation des Bases de Données}

\begin{successbox}{Optimisation des Ressources}
Une seule instance PostgreSQL et une seule MariaDB partagées par toutes les applications avec séparation logique.
\end{successbox}

\subsubsection{Architecture de Mutualisation}

\begin{figure}[H]
\centering
\begin{tikzpicture}[scale=0.8]
    % PostgreSQL Container
    \draw[fill=lightBlue, draw=primaryBlue, line width=1.5pt] (0,6) rectangle (5,8);
    \node[align=center] at (2.5,7.5) {\large\bfseries PostgreSQL};
    \node[font=\small] at (2.5,7) {thoth-postgres:5432};
    
    % Databases in PostgreSQL
    \draw[fill=codeBackground, draw=accentTeal] (0.5,6.3) rectangle (2,6.8);
    \node[font=\tiny] at (1.25,6.55) {user1\_odoo};
    
    \draw[fill=codeBackground, draw=accentTeal] (2.2,6.3) rectangle (3.7,6.8);
    \node[font=\tiny] at (2.95,6.55) {user2\_mattermost};
    
    % MariaDB Container
    \draw[fill=lightGreen, draw=secondaryGreen, line width=1.5pt] (6,6) rectangle (11,8);
    \node[align=center] at (8.5,7.5) {\large\bfseries MariaDB};
    \node[font=\small] at (8.5,7) {thoth-mariadb:3306};
    
    % Databases in MariaDB
    \draw[fill=codeBackground, draw=accentTeal] (6.5,6.3) rectangle (8,6.8);
    \node[font=\tiny] at (7.25,6.55) {user1\_wordpress};
    
    \draw[fill=codeBackground, draw=accentTeal] (8.2,6.3) rectangle (9.7,6.8);
    \node[font=\tiny] at (8.95,6.55) {user3\_moodle};
    
    % Applications
    \draw[fill=lightBlue, draw=primaryBlue] (0,3.5) rectangle (2.5,5);
    \node[font=\small] at (1.25,4.5) {Odoo};
    \node[font=\tiny] at (1.25,4.2) {user1};
    
    \draw[fill=lightBlue, draw=primaryBlue] (3,3.5) rectangle (5.5,5);
    \node[font=\small] at (4.25,4.5) {Mattermost};
    \node[font=\tiny] at (4.25,4.2) {user2};
    
    \draw[fill=lightGreen, draw=secondaryGreen] (6,3.5) rectangle (8.5,5);
    \node[font=\small] at (7.25,4.5) {WordPress};
    \node[font=\tiny] at (7.25,4.2) {user1};
    
    \draw[fill=lightGreen, draw=secondaryGreen] (9,3.5) rectangle (11.5,5);
    \node[font=\small] at (10.25,4.5) {Moodle};
    \node[font=\tiny] at (10.25,4.2) {user3};
    
    % Arrows
    \draw[->, line width=1pt, color=primaryBlue] (1.25,5) -- (1.25,6);
    \draw[->, line width=1pt, color=primaryBlue] (4.25,5) -- (2.95,6);
    \draw[->, line width=1pt, color=secondaryGreen] (7.25,5) -- (7.25,6);
    \draw[->, line width=1pt, color=secondaryGreen] (10.25,5) -- (8.95,6);
\end{tikzpicture}
\caption{Architecture de Mutualisation des Bases de Données}
\end{figure}

\subsubsection{Création de Base de Données}

\begin{lstlisting}[language=C++, caption=Provisionnement de base de données]
json PaaSOperations::provisionDatabase(const std::string& username,
                                      const std::string& appName,
                                      const std::string& dbType) {
    json result;
    
    std::string dbName = username + "_" + appName;
    std::string dbUser = dbName + "_user";
    std::string dbPass = generateRandomPassword();
    
    std::string script = "/var/lib/thoth-paas/manage-databases.sh";
    std::string cmd;
    
    if (dbType == "postgresql") {
        cmd = script + " create-pg " + username + " " + appName;
    } else if (dbType == "mysql") {
        cmd = script + " create-mysql " + username + " " + appName;
    }
    
    auto execResult = remoteExec->execute(cmd);
    
    if (execResult.success()) {
        result["success"] = true;
        result["database"] = {
            {"host", dbType == "postgresql" ? 
                    "thoth-postgres" : "thoth-mariadb"},
            {"port", dbType == "postgresql" ? 5432 : 3306},
            {"database", dbName},
            {"username", dbUser},
            {"password", dbPass}
        };
    }
    
    return result;
}
\end{lstlisting}

\subsubsection{Script de Gestion des Bases}

Extrait de \texttt{manage-databases.sh}:

\begin{lstlisting}[language=bash, caption=Création base PostgreSQL]
create-pg)
    USERNAME=$2
    APPNAME=$3
    DB_NAME="${USERNAME}_${APPNAME}"
    DB_USER="${DB_NAME}_user"
    DB_PASS=$(openssl rand -base64 16)
    
    # Create database and user
    docker exec thoth-postgres psql -U postgres <<EOSQL
        CREATE DATABASE ${DB_NAME};
        CREATE USER ${DB_USER} WITH PASSWORD '${DB_PASS}';
        GRANT ALL PRIVILEGES ON DATABASE ${DB_NAME} TO ${DB_USER};
EOSQL
    
    echo "$DB_NAME|$DB_USER|$DB_PASS"
    ;;
\end{lstlisting}

\subsection{Déploiement d'Applications}

\subsubsection{Exemple: WordPress}

\begin{lstlisting}[language=C++, caption=Déploiement WordPress]
json PaaSOperations::deployWordPress(const std::string& username,
                                    const std::string& appName) {
    json result;
    
    // 1. Provision database
    json dbResult = provisionDatabase(username, appName, "mysql");
    
    // 2. Generate docker-compose
    std::string composeFile = generateWordPressCompose(
        username, appName, dbResult["database"]);
    
    // 3. Deploy to Swarm
    std::string stackName = username + "_" + appName;
    std::string deployCmd = 
        "docker stack deploy -c " + composeFile + " " + stackName;
    
    auto deployResult = remoteExec->execute(deployCmd);
    
    result["success"] = deployResult.success();
    result["url"] = "http://" + getNodeIP() + ":" + 
                   std::to_string(assignedPort);
    
    return result;
}
\end{lstlisting}

\begin{lstlisting}[language=yaml, caption=Docker Compose WordPress]
version: '3.8'
services:
  wordpress:
    image: wordpress:latest
    environment:
      WORDPRESS_DB_HOST: thoth-mariadb
      WORDPRESS_DB_NAME: ${DB_NAME}
      WORDPRESS_DB_USER: ${DB_USER}
      WORDPRESS_DB_PASSWORD: ${DB_PASS}
    ports:
      - "${PORT}:80"
    networks:
      - paas_network
    deploy:
      replicas: 1
      
networks:
  paas_network:
    external: true
    name: thoth_paas_network
\end{lstlisting}

\subsection{Système de Facturation}

\begin{infobox}{Modèle de Facturation}
Facturation basée sur la consommation réelle des ressources (CPU, RAM, disque, durée).
\end{infobox}

\subsubsection{Calcul des Coûts}

\begin{lstlisting}[language=C++, caption=Calcul de facturation]
struct BillingInfo {
    double cpuHours;        // vCPU-heures
    double ramGBHours;      // GB-heures
    double diskGB;          // GB utilisés
    double durationHours;   // Durée totale
    
    double calculateCost() const {
        const double CPU_RATE = 10.0;     // FCFA/vCPU/heure
        const double RAM_RATE = 5.0;      // FCFA/GB/heure
        const double DISK_RATE = 0.5;     // FCFA/GB/mois
        
        double cpuCost = cpuHours * CPU_RATE;
        double ramCost = ramGBHours * RAM_RATE;
        double diskCost = diskGB * DISK_RATE * 
                         (durationHours / 720.0); // 720h = 1 mois
        
        return cpuCost + ramCost + diskCost;
    }
};

json getUserBilling(const std::string& username) {
    json result;
    double totalCost = 0.0;
    
    // Get user's applications
    auto apps = listUserApplications(username);
    
    for (const auto& app : apps) {
        BillingInfo billing = calculateAppBilling(app);
        totalCost += billing.calculateCost();
    }
    
    result["username"] = username;
    result["totalCost"] = totalCost;
    result["currency"] = "FCFA";
    
    return result;
}
\end{lstlisting}

\subsubsection{Exemple de Facturation}

\begin{table}[H]
\centering
\begin{tabular}{@{}lrrr@{}}
\toprule
\textbf{Ressource} & \textbf{Consommation} & \textbf{Tarif} & \textbf{Coût} \\
\midrule
vCPU & 2 vCPU × 720h & 10 FCFA/vCPU/h & 14,400 FCFA \\
RAM & 4 GB × 720h & 5 FCFA/GB/h & 14,400 FCFA \\
Disque & 20 GB × 1 mois & 0.5 FCFA/GB/mois & 10 FCFA \\
\midrule
\textbf{Total} & & & \textbf{28,810 FCFA} \\
\bottomrule
\end{tabular}
\caption{Exemple de Facturation Mensuelle}
\end{table}

% Partie 4
\section{Partie 4: Service SaaS}

\subsection{Plateforme Collaborative}

\begin{successbox}{Google Workspace Alternative}
OwnCloud + OnlyOffice fournissent une suite complète de collaboration: stockage, partage, édition de documents.
\end{successbox}

\subsubsection{Architecture SaaS}

\begin{figure}[H]
\centering
\begin{tikzpicture}[scale=0.85]
    % User Layer
    \draw[fill=lightBlue, draw=primaryBlue, line width=1.5pt] (0,8) rectangle (12,9.5);
    \node[align=center] at (6,8.75) {\large\bfseries Utilisateurs Web (Navigateur)};
    
    % OwnCloud
    \draw[fill=lightGreen, draw=secondaryGreen, line width=1.5pt] (0,5.5) rectangle (5.5,7.5);
    \node[align=center] at (2.75,7) {\large\bfseries OwnCloud};
    \node[font=\small] at (2.75,6.5) {Stockage \& Partage};
    \node[font=\small] at (2.75,6) {Port: 8080};
    
    % OnlyOffice
    \draw[fill=lightBlue, draw=primaryBlue, line width=1.5pt] (6.5,5.5) rectangle (12,7.5);
    \node[align=center] at (9.25,7) {\large\bfseries OnlyOffice};
    \node[font=\small] at (9.25,6.5) {Édition Documents};
    \node[font=\small] at (9.25,6) {Port: 8081};
    
    % Database Layer
    \draw[fill=codeBackground, draw=accentTeal, line width=1pt] (1,3) rectangle (4,4.5);
    \node[font=\small\bfseries] at (2.5,4) {MySQL/MariaDB};
    \node[font=\tiny] at (2.5,3.5) {Port: 3306};
    
    \draw[fill=codeBackground, draw=accentTeal, line width=1pt] (5,3) rectangle (8,4.5);
    \node[font=\small\bfseries] at (6.5,4) {Redis Cache};
    \node[font=\tiny] at (6.5,3.5) {Port: 6379};
    
    % Storage
    \draw[fill=lightGreen, draw=secondaryGreen, line width=1pt] (9,3) rectangle (11.5,4.5);
    \node[font=\small\bfseries] at (10.25,4) {Volumes};
    \node[font=\tiny] at (10.25,3.5) {Persistent};
    
    % Arrows
    \draw[->, line width=1.5pt, color=primaryBlue] (6,8) -- (2.75,7.5);
    \draw[->, line width=1.5pt, color=primaryBlue] (6,8) -- (9.25,7.5);
    \draw[->, line width=1pt, color=secondaryGreen] (2.75,5.5) -- (2.5,4.5);
    \draw[->, line width=1pt, color=secondaryGreen] (2.75,5.5) -- (6.5,4.5);
    \draw[->, line width=1pt, color=accentTeal] (5,6.5) -- (6.5,6.5);
\end{tikzpicture}
\caption{Architecture SaaS - OwnCloud + OnlyOffice}
\end{figure}

\subsubsection{Déploiement Docker Compose}

Extrait de \texttt{ownCloudSetup.sh}:

\begin{lstlisting}[language=yaml, caption=Docker Compose SaaS]
version: '3.8'

services:
  mysql:
    image: mariadb:11
    container_name: thoth-workspace-mysql
    environment:
      MYSQL_ROOT_PASSWORD: ${MYSQL_ROOT_PASS}
      MYSQL_DATABASE: owncloud
      MYSQL_USER: owncloud
      MYSQL_PASSWORD: ${MYSQL_OWNCLOUD_PASS}
    volumes:
      - mysql_data:/var/lib/mysql
    networks:
      - workspace_network

  redis:
    image: redis:7-alpine
    container_name: thoth-workspace-redis
    volumes:
      - redis_data:/data
    networks:
      - workspace_network

  owncloud:
    image: owncloud/server:latest
    container_name: thoth-owncloud
    depends_on:
      - mysql
      - redis
    environment:
      OWNCLOUD_DOMAIN: cloud.thoth.local
      OWNCLOUD_DB_TYPE: mysql
      OWNCLOUD_DB_HOST: mysql
      OWNCLOUD_ADMIN_USERNAME: admin
      OWNCLOUD_ADMIN_PASSWORD: ${ADMIN_PASS}
      OWNCLOUD_REDIS_ENABLED: "true"
      OWNCLOUD_REDIS_HOST: redis
    volumes:
      - owncloud_data:/mnt/data
    ports:
      - "8080:8080"
    networks:
      - workspace_network

  onlyoffice:
    image: onlyoffice/documentserver:latest
    container_name: thoth-onlyoffice
    environment:
      JWT_ENABLED: "true"
      JWT_SECRET: "thoth_jwt_secret_2024"
    volumes:
      - onlyoffice_data:/var/www/onlyoffice/Data
    ports:
      - "8081:80"
    networks:
      - workspace_network

networks:
  workspace_network:
    driver: bridge
    name: thoth_workspace_network
\end{lstlisting}

\subsubsection{Intégration OwnCloud + OnlyOffice}

\begin{lstlisting}[language=bash, caption=Configuration de l'intégration]
# Install OnlyOffice app in OwnCloud
docker exec thoth-owncloud occ market:install onlyoffice
docker exec thoth-owncloud occ app:enable onlyoffice

# Configure OnlyOffice integration
docker exec thoth-owncloud occ config:system:set \
    onlyoffice DocumentServerUrl \
    --value="http://thoth-onlyoffice/"

docker exec thoth-owncloud occ config:system:set \
    onlyoffice jwt_secret \
    --value="thoth_jwt_secret_2024"
\end{lstlisting}

\subsection{Provisionnement d'Utilisateurs}

\begin{lstlisting}[language=bash, caption=Script de provisionnement]
#!/bin/bash
# provision-user.sh

USERNAME=$1
PASSWORD=$2
EMAIL=$3
QUOTA=${4:-10GB}

echo "Provisioning user: $USERNAME"

# Create user in OwnCloud
docker exec thoth-owncloud occ user:add "$USERNAME" \
    --password-from-env \
    --email="$EMAIL" <<< "$PASSWORD"

# Set storage quota
docker exec thoth-owncloud occ user:setting "$USERNAME" \
    files quota "$QUOTA"

echo "User $USERNAME created with $QUOTA quota"
echo "Access: http://cloud.thoth.local:8080"
\end{lstlisting}

\subsection{Facturation Stockage}

\begin{lstlisting}[language=bash, caption=Calcul facturation stockage]
#!/bin/bash
# billing-tracker.sh

OWNCLOUD_DATA="/var/lib/thoth-saas/owncloud"
BILLING_RATE=0.10  # $ per GB per month

calculate_storage() {
    USAGE_BYTES=$(du -sb "$OWNCLOUD_DATA" | cut -f1)
    USAGE_GB=$(echo "scale=2; $USAGE_BYTES/1024/1024/1024" | bc)
    MONTHLY_COST=$(echo "scale=2; $USAGE_GB * $BILLING_RATE" | bc)
    
    cat > billing.json <<JSON
{
  "storage_gb": ${USAGE_GB},
  "rate_per_gb": ${BILLING_RATE},
  "monthly_cost": ${MONTHLY_COST},
  "timestamp": "$(date -Iseconds)"
}
JSON
    
    echo "Storage: ${USAGE_GB} GB"
    echo "Cost: \$${MONTHLY_COST}/month"
}
\end{lstlisting}

\begin{table}[H]
\centering
\begin{tabular}{@{}lrrr@{}}
\toprule
\textbf{Utilisateur} & \textbf{Stockage} & \textbf{Tarif} & \textbf{Coût} \\
\midrule
alice & 5.2 GB & 0.10 \$/GB & 0.52 \$ \\
bob & 12.8 GB & 0.10 \$/GB & 1.28 \$ \\
charlie & 3.1 GB & 0.10 \$/GB & 0.31 \$ \\
\midrule
\textbf{Total} & 21.1 GB & & \textbf{2.11 \$} \\
\bottomrule
\end{tabular}
\caption{Exemple de Facturation Stockage}
\end{table}

% API Documentation
\section{Documentation API}

\subsection{Authentification}

\begin{infobox}{Token-Based Authentication}
L'API utilise des tokens de session avec expiration d'1 heure.
\end{infobox}

\subsubsection{Login}

\begin{lstlisting}[language=bash, caption=Endpoint de connexion]
POST /api/auth/login
Content-Type: application/json

{
  "username": "alice",
  "password": "secret123"
}

Response:
{
  "success": true,
  "token": "a1b2c3d4e5f6...",
  "user": {
    "username": "alice",
    "role": "user",
    "email": "alice@example.com"
  }
}
\end{lstlisting}

\subsection{Gestion des VMs}

\subsubsection{Déployer une VM}

\begin{lstlisting}[language=bash, caption=Création de VM]
POST /api/vms/deploy
Authorization: Bearer {token}
Content-Type: application/json

{
  "hostname": "webserver",
  "flavor": "medium",
  "baseImage": "ubuntu-22.04",
  "network": "user_network_id",
  "username": "ubuntu",
  "password": "mypass123"
}

Response:
{
  "success": true,
  "vm": {
    "name": "alice__webserver__1234567890",
    "status": "running",
    "ip": "192.168.100.50"
  }
}
\end{lstlisting}

\subsubsection{Lister les VMs}

\begin{lstlisting}[language=bash, caption=Liste des VMs utilisateur]
GET /api/vms/list
Authorization: Bearer {token}

Response:
{
  "success": true,
  "vms": [
    {
      "name": "alice__webserver__1234567890",
      "displayName": "webserver",
      "status": "running",
      "ip": "192.168.100.50",
      "vcpus": 2,
      "memory": 4096,
      "created": 1704067200000
    }
  ]
}
\end{lstlisting}

\subsection{Gestion Swarm}

\subsubsection{Créer un Cluster}

\begin{lstlisting}[language=bash, caption=Création cluster Swarm]
POST /api/swarm/create
Authorization: Bearer {token}
Content-Type: application/json

{
  "clusterName": "production",
  "numManagers": 1,
  "numWorkers": 2
}

Response:
{
  "success": true,
  "cluster": {
    "clusterId": "cluster_abc123",
    "clusterName": "production",
    "nodes": [
      {
        "name": "production-manager-1",
        "role": "manager",
        "ip": "192.168.100.101"
      },
      {
        "name": "production-worker-1",
        "role": "worker",
        "ip": "192.168.100.102"
      }
    ]
  }
}
\end{lstlisting}

\subsection{Déploiement PaaS}

\begin{lstlisting}[language=bash, caption=Déploiement application]
POST /api/paas/deploy
Authorization: Bearer {token}
Content-Type: application/json

{
  "appType": "wordpress",
  "appName": "myblog",
  "cluster": "production"
}

Response:
{
  "success": true,
  "application": {
    "appId": "app_xyz789",
    "appName": "myblog",
    "type": "wordpress",
    "url": "http://192.168.100.101:8001",
    "database": {
      "host": "thoth-mariadb",
      "database": "alice_myblog",
      "username": "alice_myblog_user"
    }
  }
}
\end{lstlisting}

% Deployment Guide
\section{Guide de Déploiement}

\subsection{Prérequis}

\begin{enumerate}
    \item \textbf{Système}: Ubuntu 22.04 LTS ou supérieur
    \item \textbf{CPU}: Support de virtualisation (Intel VT-x / AMD-V)
    \item \textbf{RAM}: Minimum 16 GB recommandé
    \item \textbf{Disque}: 100 GB minimum
    \item \textbf{Réseau}: Connexion Internet
\end{enumerate}

\subsection{Installation}

\subsubsection{1. Cloner le Repository}

\begin{lstlisting}[language=bash, caption=Clonage du projet]
git clone https://github.com/votre-org/thoth-cloud.git
cd thoth-cloud
\end{lstlisting}

\subsubsection{2. Setup KVM Host}

\begin{lstlisting}[language=bash, caption=Configuration hôte KVM]
# Run setup script as root
sudo bash scripts/setup_remote.sh

# This installs:
# - KVM, QEMU, libvirt
# - Creates 'vps' user
# - Configures passwordless sudo
# - Sets up firewall rules
\end{lstlisting}

\subsubsection{3. Compiler le Backend}

\begin{lstlisting}[language=bash, caption=Compilation C++]
# Install dependencies
sudo bash scripts/setup-libs.sh

# Create build directory
mkdir build && cd build

# Compile
cmake ..
make -j$(nproc)

# Binary will be at: build/thoth_cloud
\end{lstlisting}

\subsubsection{4. Configurer l'Environnement}

\begin{lstlisting}[language=bash, caption=Configuration .env]
# Create config directory
sudo mkdir -p /var/lib/thoth-cloud

# Create .env file
cat > /var/lib/thoth-cloud/.env <<EOF
API_PORT=3000
HOSTS=qemu:///system
EOF
\end{lstlisting}

\subsubsection{5. Initialiser les Bases de Données}

\begin{lstlisting}[language=bash, caption=Setup bases partagées]
sudo bash scripts/paas-db-setup.sh

# This creates:
# - thoth-postgres container
# - thoth-mariadb container
# - Database management scripts
\end{lstlisting}

\subsubsection{6. Démarrer le Backend}

\begin{lstlisting}[language=bash, caption=Lancement du serveur]
# Run backend
sudo ./build/thoth_cloud

# Server starts on http://localhost:3000
# Check logs in /var/log/thoth-cloud/
\end{lstlisting}

\subsubsection{7. Créer un Utilisateur Admin}

\begin{lstlisting}[language=bash, caption=Création admin]
# Compile and run admin creator
g++ -o create_admin src/create_admin.cpp \
    -lvirt -lcrypto $(pkg-config --cflags --libs)

sudo ./create_admin

# Default admin credentials:
# Username: admin
# Password: admin123
\end{lstlisting}

\subsection{Multi-Hôtes (Optionnel)}

\begin{lstlisting}[language=bash, caption=Ajout hôte distant]
# On remote host, run:
sudo bash scripts/set-remote.sh

# Get remote host IP
REMOTE_IP=$(hostname -I | awk '{print $1}')

# On backend server, generate SSH key
ssh-keygen -t rsa -b 4096 -f ~/.ssh/thoth_kvm_key -N ''

# Copy key to remote
ssh-copy-id -i ~/.ssh/thoth_kvm_key.pub vps@$REMOTE_IP

# Test connection
virsh -c qemu+ssh://vps@$REMOTE_IP/system list

# Add to configuration
# Edit /var/lib/thoth-cloud/.env:
# HOSTS=qemu:///system,qemu+ssh://vps@$REMOTE_IP/system
\end{lstlisting}

\subsection{Vérification}

\begin{lstlisting}[language=bash, caption=Tests de vérification]
# Test API health
curl http://localhost:3000/api/health

# Login
curl -X POST http://localhost:3000/api/auth/login \
  -H "Content-Type: application/json" \
  -d '{"username":"admin","password":"admin123"}'

# List VMs
TOKEN="your_token_here"
curl http://localhost:3000/api/vms/list \
  -H "Authorization: Bearer $TOKEN"
\end{lstlisting}

% Limitations and Improvements
\section{Limites et Améliorations}

\subsection{Limites Connues}

\begin{enumerate}
    \item \textbf{Sécurité}:
    \begin{itemize}
        \item Tokens stockés en mémoire (perdus au redémarrage)
        \item Pas de rotation automatique des secrets
        \item Authentification basique (pas de 2FA)
    \end{itemize}
    
    \item \textbf{Scalabilité}:
    \begin{itemize}
        \item Backend mono-instance
        \item Pas de load balancing natif
        \item Limite de 20 hôtes KVM
    \end{itemize}
    
    \item \textbf{Monitoring}:
    \begin{itemize}
        \item Métriques limitées
        \item Pas d'alerting automatique
        \item Logs non centralisés
    \end{itemize}
    
    \item \textbf{Réseau}:
    \begin{itemize}
        \item Gestion basique des sous-réseaux
        \item Pas de VLAN support
        \item Firewall rules manuelles
    \end{itemize}
\end{enumerate}

\subsection{Améliorations Futures}

\begin{successbox}{Roadmap}
Améliorations prévues pour les prochaines versions
\end{successbox}

\subsubsection{Court Terme (3 mois)}

\begin{enumerate}
    \item \textbf{Authentification}:
    \begin{itemize}
        \item OAuth2/OpenID Connect
        \item LDAP/Active Directory integration
        \item Two-factor authentication (2FA)
        \item Session persistence (Redis)
    \end{itemize}
    
    \item \textbf{Monitoring}:
    \begin{itemize}
        \item Intégration Prometheus/Grafana
        \item Métriques détaillées (CPU, RAM, réseau, I/O)
        \item Alerting (email, Slack, webhook)
        \item Dashboard temps réel
    \end{itemize}
    
    \item \textbf{Backups}:
    \begin{itemize}
        \item Snapshots automatiques
        \item Backup incrémental
        \item Restauration rapide
        \item Réplication cross-host
    \end{itemize}
\end{enumerate}

\subsubsection{Moyen Terme (6 mois)}

\begin{enumerate}
    \item \textbf{HA et Clustering}:
    \begin{itemize}
        \item Backend multi-instances
        \item Load balancer (HAProxy/Nginx)
        \item Failover automatique
        \item Shared storage (Ceph/GlusterFS)
    \end{itemize}
    
    \item \textbf{Networking Avancé}:
    \begin{itemize}
        \item SDN (Software-Defined Networking)
        \item VLAN support
        \item Firewall programmable
        \item VPN inter-sites
    \end{itemize}
    
    \item \textbf{Kubernetes}:
    \begin{itemize}
        \item Support K8s en plus de Swarm
        \item Helm charts
        \item Operators personnalisés
        \item Multi-cluster management
    \end{itemize}
\end{enumerate}

\subsubsection{Long Terme (12 mois)}

\begin{enumerate}
    \item \textbf{AI/ML Integration}:
    \begin{itemize}
        \item Auto-scaling prédictif
        \item Anomaly detection
        \item Resource optimization
        \item Cost forecasting
    \end{itemize}
    
    \item \textbf{Marketplace}:
    \begin{itemize}
        \item Catalogue d'images personnalisées
        \item Templates d'infrastructure
        \item Partage communautaire
        \item Validation automatique
    \end{itemize}
    
    \item \textbf{Compliance}:
    \begin{itemize}
        \item Audit trails
        \item GDPR compliance tools
        \item Encryption at rest
        \item Security scanning
    \end{itemize}
\end{enumerate}

% Conclusion
\section{Conclusion}

\begin{successbox}{Réalisations}
THOTH CLOUD implémente avec succès une plateforme cloud complète intégrant IaaS, PaaS et SaaS.
\end{successbox}

\subsection{Objectifs Atteints}

\begin{enumerate}
    \item \textbf{IaaS Multi-Tenant}:
    \begin{itemize}
        \item Isolation utilisateurs complète
        \item Gestion multi-hôtes avec 3 stratégies
        \item Système de flavors avec pricing
        \item Réseaux dédiés par utilisateur
    \end{itemize}
    
    \item \textbf{PaaS Docker Swarm}:
    \begin{itemize}
        \item Déploiement automatique de clusters
        \item 8 applications en catalogue
        \item Bases de données mutualisées
        \item Facturation à la consommation
    \end{itemize}
    
    \item \textbf{SaaS Collaboratif}:
    \begin{itemize}
        \item OwnCloud + OnlyOffice intégrés
        \item Alternative Google Workspace
        \item Facturation au stockage
        \item Multi-utilisateurs
    \end{itemize}
\end{enumerate}

\subsection{Compétences Développées}

\begin{itemize}
    \item Architecture cloud multi-services
    \item Programmation système en C++
    \item Virtualisation KVM/libvirt avancée
    \item Orchestration de conteneurs
    \item Gestion de réseaux virtuels
    \item Développement d'API REST
    \item DevOps et automatisation
\end{itemize}

\subsection{Perspectives}

THOTH CLOUD constitue une base solide pour un cloud de production. Les améliorations futures permettront d'atteindre la maturité d'un cloud professionnel tout en conservant la simplicité d'utilisation.

\vspace{1cm}

\begin{center}
\begin{tikzpicture}
    \draw[fill=lightGreen, draw=secondaryGreen, line width=2pt] (0,0) rectangle (12,2);
    \node[align=center, font=\Large\bfseries] at (6,1.3) {\color{darkBlue}THOTH CLOUD};
    \node[align=center, font=\large] at (6,0.6) {\color{primaryBlue}Cloud Computing for Everyone};
\end{tikzpicture}
\end{center}

% Appendices
\appendix

\section{Structure du Code}

\subsection{Organisation des Fichiers}

\begin{lstlisting}[language=bash, caption=Arborescence du projet]
thoth-cloud/
├── include/               # Header files (.hpp)
│   ├── baseimage_manager.hpp
│   ├── config_manager.hpp
│   ├── cors.hpp
│   ├── flavor_manager.hpp
│   ├── host_manager.hpp
│   ├── network_manager.hpp
│   ├── paas_operations.hpp
│   ├── remote_executor.hpp
│   ├── routes.hpp
│   ├── swarm_operations.hpp
│   ├── user_operations.hpp
│   ├── utils.hpp
│   ├── validation.hpp
│   └── vm_operations.hpp
├── src/                   # Implementation files (.cpp)
│   ├── baseimage_manager.cpp
│   ├── config_manager.cpp
│   ├── cors.cpp
│   ├── flavor_manager.cpp
│   ├── host_manager.cpp
│   ├── main.cpp
│   ├── network_manager.cpp
│   ├── paas_operations.cpp
│   ├── remote_executor.cpp
│   ├── resource_logger.cpp
│   ├── routes.cpp
│   ├── routes2.cpp
│   ├── swarm_operations.cpp
│   ├── user_operations.cpp
│   ├── utils.cpp
│   ├── validation.cpp
│   └── vm_operations.cpp
├── scripts/               # Deployment scripts
│   ├── general.sh
│   ├── setup_remote.sh
│   ├── paas-db-setup.sh
│   ├── ownCloudSetup.sh
│   └── setup-novnc.sh
└── build/                 # Compiled binaries
\end{lstlisting}

\section{Configuration Système}

\subsection{Fichiers de Configuration}

\subsubsection{/var/lib/thoth-cloud/.env}

\begin{lstlisting}[caption=Fichier d'environnement]
API_PORT=3000
HOSTS=qemu:///system,qemu+ssh://vps@192.168.1.100/system
BASE_IMAGE_DIR=/var/lib/libvirt/images/baseimg
CLOUD_INIT_DIR=/var/lib/libvirt/images/cloud-init-iso
\end{lstlisting}

\subsubsection{/var/lib/thoth-cloud/hosts.json}

\begin{lstlisting}[language=json, caption=Configuration des hôtes]
{
  "hosts": [
    {
      "id": "localhost",
      "uri": "qemu:///system",
      "hostname": "localhost"
    },
    {
      "id": "kvm-node-1",
      "uri": "qemu+ssh://vps@192.168.1.100/system",
      "hostname": "kvm-node-1"
    }
  ]
}
\end{lstlisting}

\section{Commandes Utiles}

\subsection{Gestion des VMs}

\begin{lstlisting}[language=bash, caption=Commandes virsh]
# List all VMs
virsh list --all

# Start VM
virsh start vm_name

# Shutdown VM
virsh shutdown vm_name

# Force stop
virsh destroy vm_name

# Delete VM
virsh undefine vm_name --remove-all-storage

# Get VM info
virsh dominfo vm_name

# Console access
virsh console vm_name
\end{lstlisting}

\subsection{Gestion Docker Swarm}

\begin{lstlisting}[language=bash, caption=Commandes Swarm]
# Initialize swarm
docker swarm init --advertise-addr <IP>

# Join as worker
docker swarm join --token <TOKEN> <MANAGER_IP>:2377

# List nodes
docker node ls

# Deploy stack
docker stack deploy -c compose.yml stack_name

# List services
docker service ls

# Scale service
docker service scale service_name=5

# Remove stack
docker stack rm stack_name
\end{lstlisting}

\end{document}